\section{The prime-degree operator and its exact Weil pairing}
\label{sec:relative-prime-degree}
We now use the relative principal-parts object produced by the full
spectral extension, rather than only comparing it with older determinant
bounds. At $\Spec\Z$ its arithmetic part is $\Q/\Z$, and its $p$-primary
part is
\[
 D_p=\Z[1/p]/\Z\xrightarrow{\ \sim\ }\Q_p/\Z_p.
 \tag{PD1}
\]
The relative cohomological construction and its receiving category are
proved in the preceding recollement chapter. The operators and pairings
on (PD1) will all be calculated here. The absolute value remains
$|p|_p=p^{-1}$, and the additive Haar measure remains $\mu(\Z_p)=1$.
The use of local Fourier transforms follows the same convention as
Connes \cite[Section III]{Connes1998}. No spectral interpretation of the
zeros is assumed.

\subsection{The actual filtration, degree and finite character pairing}
Put $D_{p,n}=p^{-n}\Z/\Z$, for $n\ge0$. This is a subgroup of
$D_p$ of order $p^n$, and $D_p=\bigcup_nD_{p,n}$. Mapping rational
classes into $\Q_p/\Z_p$ is injective: an element $a/p^n$ lying in
$\Z_p$ has $p^n\mid a$ after cancellation. It is onto: the finite
negative part of any $p$-adic expansion differs from the number by an
element of $\Z_p$. This proves (PD1), including its filtration.
Multiplication by $p$ gives
\[
 0\longrightarrow D_{p,1}\longrightarrow D_p
 \xrightarrow{p}D_p\longrightarrow0,
 \qquad |D_{p,1}|=p.
 \tag{PD2}
\]
Surjectivity follows by replacing $a/p^n$ by $a/p^{n+1}$; the kernel
consists of the $p$ classes $a/p$ modulo integers. Every fibre therefore
has exactly $p$ elements. The exact layer maps are
\[
 D_{p,n}/D_{p,n-1}\xrightarrow{\sim}\F_p,
 \qquad [a/p^n]\longmapsto a\pmod p,\quad n\ge1.
 \tag{PD3}
\]
They are well defined because changing the representative modulo
$D_{p,n-1}$ changes $a$ by a multiple of $p$. They are injective and
onto by the same formula. Thus each layer has module length one and
logarithmic cardinality $\log p$.

For $d=a/p^n\pmod\Z$ and $x\in\Z_p$, define
\[
 \chi_d(x)=\exp\left(2\pi i\,\frac{a(x\bmod p^n)}{p^n}\right).
 \tag{PD4}
\]
Changing representatives of either residue changes the exponent by an
integer; passing to denominator $p^{n+1}$ gives the same character.
At finite level this is a perfect pairing of $D_{p,n}$ and
$\Z/p^n\Z$: a nonzero $a$ is detected by $x=1$, and a nonzero $x$
is detected by $a=1$. The geometric-series identity gives
\[
 p^{-n}\sum_{x\bmod p^n}\chi_d(x)\overline{\chi_{d'}(x)}
 =\mathbf1_{d=d'}\quad(d,d'\in D_{p,n}).
 \tag{PD5}
\]
The Haar measure of each residue class is $p^{-n}$, so these are also
the integrals over $\Z_p$.

Let $\ell^2(D_p)$ use counting measure. Formula (PD5) implies that
\[
 \mathcal F_p:\ell^2(D_p)\longrightarrow L^2(\Z_p,dx),\qquad
 \mathcal F_p\delta_d=\chi_d
 \tag{PD6}
\]
extends isometrically from finite sums. It is onto: the $p^n$ orthogonal
characters at level $n$ form a basis of the $p^n$-dimensional space of
functions constant on residue classes modulo $p^n$. The union of these
spaces is dense in $L^2(\Z_p)$. One way to see this last assertion is
to approximate a continuous function uniformly using uniform continuity
on the compact space $\Z_p$ and then use the density of continuous
functions for the finite regular Haar measure. In particular
$\mathcal F_p\delta_0=\mathbf1_{\Z_p}$, with norm exactly one.

\subsection{The square-root factor is fixed by the prime degree}
The pullback of the actual map (PD2) on finite-support functions is
\[
 (A_pf)(d)=f(pd),\qquad d\in D_p.
 \tag{PD7}
\]
Each fibre has $p$ elements, so direct summation proves
$\|A_pf\|^2=p\|f\|^2$. Consequently its bounded extension has
norm exactly $\sqrt p$. Define
$V_p=p^{-1/2}A_p$ and keep both operators distinguished. The factor
in $V_p$ is forced by this degree calculation.
Their adjoints and the Fourier transforms are
\begin{align}
 (V_p^*f)(u)&=p^{-1/2}\sum_{pd=u}f(d),\tag{PD8}\\
 (\mathcal F_pV_p\mathcal F_p^{-1}F)(x)
 &=\sqrt p\,\mathbf1_{p\Z_p}(x)F(x/p),\tag{PD9}\\
 (\mathcal F_pV_p^*\mathcal F_p^{-1}F)(x)
 &=p^{-1/2}F(px).\tag{PD10}
\end{align}
\begin{proof}
Taking the counting-measure inner product in (PD7) and collecting terms
over each fibre proves (PD8). For finite $f$, the Fourier transform of
$A_pf$ is
$\sum_u f(u)\sum_{pd=u}\chi_d(x)$. The inner character sum is zero
unless $x\in p\Z_p$: translation of the fibre by $D_{p,1}$ multiplies
it by a nontrivial $p$th root of unity otherwise. For $x=py$ it is
$p\chi_u(y)$. This proves (PD9) on the dense finite-support subspace
and hence on the whole Hilbert space. Finally changing variables
$x=py$, with $dx=p^{-1}dy$, in the inner product for (PD9) proves
(PD10). No counting measure or Haar factor has been omitted.
\end{proof}

In particular $V_p^*V_p=I$. In the Fourier representation
$V_p^n(V_p^*)^n$ is multiplication by $\mathbf1_{p^n\Z_p}$.
Their decreasing intersection has measure zero, since
$\mu(p^n\Z_p)=p^{-n}$. Thus $\bigcap_nV_p^n\ell^2(D_p)=0$.
The distinction between this isometry and an invertible evolution is
resolved by the following explicit extension, not by an assumed inverse.

\begin{theorem}[Exact unitary dilation and its original vector]
\label{thm:PD-dilation}
On the actual local space $L^2(\Q_p,dx)$ define
\[
 (\mathscr U_pF)(x)=\sqrt p\,F(x/p),\qquad
 (\mathscr U_p^{-1}F)(x)=p^{-1/2}F(px).
 \tag{PD11}
\]
This is unitary. Extension by zero embeds $L^2(\Z_p)$ isometrically,
and $\mathscr U_p$ restricts there to (PD9). The union of
$\mathscr U_p^{-n}L^2(\Z_p)$ is dense in $L^2(\Q_p)$.
For the original vector $\xi_p=\mathbf1_{\Z_p}$,
\[
 \langle \mathscr U_p^n\xi_p,\xi_p\rangle=p^{-|n|/2},\qquad n\in\Z.
 \tag{PD12}\label{eq:PD-moments}
\]
Its scalar spectral measure on the unit circle is
\[
 d\nu_p(\theta)=\frac1{2\pi}
 \frac{1-p^{-1}}{1-2p^{-1/2}\cos\theta+p^{-1}}\,d\theta.
 \tag{PD13}\label{eq:PD-measure}
\]
\end{theorem}
\begin{proof}
Change of variables gives
$\int p|F(x/p)|^2dx=\int|F(y)|^2dy$, and the two formulas (PD11)
compose to the identity. A function supported in $\Z_p$ is carried
to one supported in $p\Z_p$, proving the restriction. The inverse
images are exactly the spaces supported in $p^{-n}\Z_p$. These sets
exhaust $\Q_p$, so truncation of any $L^2$ function proves density.
For $n\ge0$,
$\mathscr U_p^n\xi_p=p^{n/2}\mathbf1_{p^n\Z_p}$; pairing with
$\mathbf1_{\Z_p}$ gives $p^{n/2}p^{-n}=p^{-n/2}$.
Unitarity gives the complex conjugate value for $-n$, here the same real
number. This proves (PD12).

For $0<r<1$, absolutely convergent geometric series give
$\sum_{n\in\Z}r^{|n|}e^{in\theta}
=(1-r^2)/(1-2r\cos\theta+r^2)$. Its constant coefficient is one,
and its density is positive. Setting $r=p^{-1/2}$ shows that (PD13)
has exactly the moments (PD12). On any Laurent polynomial $a$, direct
expansion then gives
$\|a(\mathscr U_p)\xi_p\|^2=\int|a(e^{i\theta})|^2d\nu_p(\theta)$.
Density of Laurent polynomials identifies the cyclic Hilbert space with
$L^2(\nu_p)$ and $\mathscr U_p$ with multiplication by $e^{i\theta}$ there.
This constructs the asserted scalar spectral measure directly.
\end{proof}

\subsection{The original prime term is the energy of this operator}
Write $a_p=\log p$ and
$f(u)=g(e^u)$ for the original multiplicative test $g$.
Define a vector-valued function on the interval $[0,a_p)$ by
\[
 (\mathcal E_pg)(x)=
 \sum_{n\in\Z} f(x+n a_p)\mathscr U_p^n\xi_p
 \ \in L^2(\Q_p),\qquad 0\le x<a_p.
 \tag{PD14}\label{eq:PD-energy-map}
\]
For compact support of $f$, only finitely many $n$ occur, uniformly in
this bounded interval. Thus this is an explicitly defined map into
$L^2([0,a_p),dx;L^2(\Q_p))$.

\begin{theorem}[The relative-prime Hilbert pairing in Weil's formula]
\label{thm:PD-Weil-energy}
For every $g\in C_c^\infty(\R_{>0})$, with $h=g*g^*$,
\begin{align}
 \|\mathcal E_pg\|^2
 &=\sum_{j\in\Z}p^{-|j|/2}h(p^j)\tag{PD15}\\
 &=\frac1{2\pi}\int_\R|\widehat g(t)|^2
 P_{p^{-1/2}}(t\log p)\,dt.\tag{PD16}
\end{align}
In particular the complete arithmetic prime term, including every prime
power, is exactly
\[
 W_p(g*g^*)=(\log p)
 \left(\|\mathcal E_pg\|^2-\|g\|_{L^2(dx/x)}^2\right).
 \tag{PD17}\label{eq:PD-prime-energy}
\]
The full finite-place return therefore becomes
\[
 D_F(g*g^*)=\mathscr B(Pg,Pg)-W_\infty(g*g^*)
 +\sum_{p\in F}(\log p)
 \left(\|g\|_{L^2(dx/x)}^2-\|\mathcal E_pg\|^2\right).
 \tag{PD18}\label{eq:PD-full-energy}
\]
\end{theorem}
\begin{proof}
Use the inner product linear in its first variable. Expanding (PD14)
and using (PD12) gives
\[
 \sum_{n,m}\int_0^{a_p}
 f(x+n a_p)\overline{f(x+m a_p)}p^{-|n-m|/2}\,dx.
\]
Set $j=n-m$ and $u=x+n a_p$. For fixed $j$, the half-open intervals
$[n a_p,(n+1)a_p)$ partition $\R$ up to measure zero, so the sum is
$\sum_jp^{-|j|/2}\int_\R f(u)\overline{f(u-j a_p)}du$.
The integral is precisely $h(p^j)$ with the original multiplicative
convolution, proving (PD15). There are only finitely many nonzero
correlations for this test. Fourier inversion of the correlations and
the absolutely convergent Poisson series prove (PD16); the dominating
sum is $(1+r)/(1-r)|\widehat g(t)|^2$ with $r=p^{-1/2}<1$,
which is integrable. The $j=0$ term of (PD15) is
$h(1)=\|g\|^2$. Multiplying the remaining terms by $\log p$
gives the original $W_p(h)$, proving (PD17). Substitution into the
unmodified explicit formula proves (PD18).
\end{proof}

\subsection{The same determinant and the exact finite support guard}
The scalar resolvent on this original vector is also explicit:
\[
 \langle (I-z\mathscr U_p)^{-1}\xi_p,\xi_p\rangle
 =\sum_{n\ge0}z^np^{-n/2}=
 \frac1{1-zp^{-1/2}},\qquad |z|<1.
 \tag{PD19}
\]
The operator series is norm convergent in this domain. The scalar series
and its last expression continue to $|z|<\sqrt p$ without asserting
operator resolvent invertibility there. Substituting
$z=p^{1/2-s}$ for $\Re s>1/2$ identifies the scalar with the local
determinant receiver $Z_p(s)$ of (ME12); the resulting scalar identity
continues to $\Re s>0$ by the displayed convergent scalar series.
This is the exact comparison between the noninvertible prime-degree
map, its unitary extension, and the mixed-face determinant.

If $\supp g\subset[\lambda^{-1},\lambda]$ and $p>\lambda^2$,
then $h(p^j)=0$ for all $j\ne0$. Formula (PD15) gives
$\|\mathcal E_pg\|^2=\|g\|^2$ exactly, so (PD18) agrees with the
full prime return once $F$ contains the primes at most $\lambda^2$.
It is essential to take the differences before enlarging $F$; neither
of the two infinite sums of positive energies is asserted to converge
separately.

\subsection{The precise local weight identity}
The unweighted prime pullback $A_p$ has the invertible extension
\[
 T_p^{\mathrm{rel}}=\sqrt p\,\mathscr U_p,
 \qquad (T_p^{\mathrm{rel}}F)(x)=pF(x/p)
 \quad\hbox{on }L^2(\Q_p).
 \tag{PD20}
\]
The superscript distinguishes this relative-prime operator from the
finite mixed-boundary matrix $T_p=D_p$ of (AB8); the connecting scalar
resolvent was proved in (PD19).

\begin{theorem}[Weight-one norm identity on the computed relative space]
\label{thm:PD-local-weight}
For the actual operator (PD20),
\[
 (T_p^{\mathrm{rel}})^*T_p^{\mathrm{rel}}
 =T_p^{\mathrm{rel}}(T_p^{\mathrm{rel}})^*=pI,
 \qquad \sigma(T_p^{\mathrm{rel}})=\{z\in\C:|z|=\sqrt p\}.
 \tag{PD21}
\]
It has no nonzero Hilbert-space eigenvectors. The spectrum assertion
concerns the entire Hilbert spectrum, not a list of zeta eigenvalues.
\end{theorem}
\begin{proof}
The two adjoint identities follow from the proved unitarity of $\mathscr U_p$.
For the exact spectrum, put
$\mathcal A_p=\Z_p^\times$ and
$H_0=L^2(\mathcal A_p)\subset L^2(\Q_p)$. The measurable disjoint
sets $p^n\mathcal A_p$, $n\in\Z$, partition $\Q_p\setminus\{0\}$.
The missing singleton has Haar measure zero. The summands
$\mathscr U_p^nH_0$ are therefore mutually orthogonal and their Hilbert direct
sum is the whole space. Choose a nonzero unit vector $h\in H_0$;
then $h_n=\mathscr U_p^nh$ is orthonormal. For $|\lambda|=1$ put
$v_N=(2N+1)^{-1/2}\sum_{n=-N}^N\lambda^{-n}h_n$.
Cancellation at the interior indices gives
$\|(\mathscr U_p-\lambda I)v_N\|^2=2/(2N+1)$, while $\|v_N\|=1$.
Hence every such $\lambda$ is in the spectrum. For $|\lambda|>1$,
the Neumann series for $-\lambda^{-1}(I-\lambda^{-1}\mathscr U_p)^{-1}$
inverts $\mathscr U_p-\lambda I$; for $|\lambda|<1$ use
$(I-\lambda \mathscr U_p^{-1})^{-1}\mathscr U_p^{-1}$. Thus there is no other spectrum.
Scaling proves (PD21).

If $\mathscr U_pF=\lambda F$, then $|\lambda|=1$ follows by taking norms
unless $F=0$. The orthogonal components in $\mathscr U_p^nH_0$ consequently
have equal norms for all $n$. Square summability forces them all to
be zero. Scaling again proves the assertion for $T_p^{\mathrm{rel}}$.
\end{proof}

For comparison, Deligne's actual purity result in Weil II, (3.3.6),
concerns the image of compactly supported cohomology in ordinary
cohomology, using the upper-weight theorem and its dual lower-weight
theorem \cite{DeligneII}. Formula (PD21) supplies a complex weight-one
norm identity on the relative object explicitly constructed here. The
remaining programme step is to calculate the corresponding image for
the full split cohomology and its trace, with the maps (ME13), (PD14)
and the complete expression (PD18). No identification of that image
with $L^2(\Q_p)$, and no full cohomological purity, is asserted by
the local calculation.

\begin{figure}[ht]
\centering
\begin{tikzpicture}[node distance=15mm,
 every node/.style={align=center},
 box/.style={draw,rounded corners,text width=112mm,inner sep=5pt}]
\node[box] (parts) {$D_p=\Q_p/\Z_p$ in the relative principal parts\\
$d\mapsto pd$: every fibre has $p$ elements};
\node[box,below=of parts] (shift) {$A_pf(d)=f(pd)$ has norm $\sqrt p$\\
$V_p=p^{-1/2}A_p$ corresponds to
$\sqrt p\,\mathbf1_{p\Z_p}(x)F(x/p)$};
\node[box,below=of shift] (whole) {$\mathscr U_pF(x)=\sqrt pF(x/p)$ on $L^2(\Q_p)$\\
$\xi_p=\mathbf1_{\Z_p}$,
$\langle \mathscr U_p^n\xi_p,\xi_p\rangle=p^{-|n|/2}$};
\node[box,below=of whole] (energy) {Original Weil prime term (PD17):\\
$W_p(g*g^*)=(\log p)(\|\mathcal E_pg\|^2-\|g\|^2)$};
\draw[-{Latex}] (parts)--node[right]{(PD7)--(PD8)}(shift);
\draw[-{Latex}] (shift)--node[right]{isometric inclusion (PD11)}(whole);
\draw[-{Latex}] (whole)--node[right]{the exact map (PD14)}(energy);
\end{tikzpicture}
\caption{Use of the relative prime data in the original arithmetic pairing.
The prime degree fixes every factor of $p^{1/2}$; the original Haar
measure fixes every norm. Theorems~\ref{thm:PD-dilation} and
\ref{thm:PD-Weil-energy} prove the arrows. The formula retains the
subtracted diagonal energy, both poles and the Archimedean contribution
in (PD18), in the conventions of \cite{Connes1998,Connes2026}.}
\end{figure}

These calculations establish a positive local Hilbert representation
constructed from the relative prime object and the exact signed map
from it into the Weil functional. The full inequality for (PD18) has
not been proved: positivity of each individual squared norm does not
assign the sign of their displayed difference. The new task supplied
by this construction is the full pairing of these relative-prime
operators with the interior and Archimedean objects, with the
subtraction and the two endpoint terms retained.
